\documentclass{ctexart}

\usepackage{packages}
\usepackage{mathstyle}
\usepackage{reportstyle}

\setinfo{理论分析}{崔正平}

\begin{document}

\maketitle

\section{问题设置与记号}
\label{sec:setup}

考虑高对比度多尺度问题, 对其采用有限差分或者有限元格式得到对称正定线性方程组 $\bm A\bm u=\bm b$. 按细点与粗点分块,
\begin{align}
    \bm A&=\begin{bNiceMatrix}
        \bm B & \bm C\\
        \bm C^\T & \bm A_{\mathrm{CC}}
    \end{bNiceMatrix}\succ \mathbf{O},\q \bm B=\bm A_{\mathrm{FF}}\succ \mathbf{O},
    \label{eq:setup-block-A}
\end{align}
其中 $n_{\mathrm{F}},n_{\mathrm{C}}$ 分别为细点和粗点的个数. 固定粗点为单位注入后, 插值算子可写成
\begin{align}
    \bm P(\bm W)&=\begin{bNiceMatrix}
        \bm W\\
        \bm I
    \end{bNiceMatrix},\q \bm W\in\R^{n_{\mathrm{F}}\times n_{\mathrm{C}}}.
    \label{eq:setup-interpolation}
\end{align}
取前、后磨光子分别为前向、后向 Gauss--Seidel 迭代, 前磨光误差传播算子记为 $\bm G$, 后磨光误差传播算子取其 $\bm A$-伴随, Galerkin 粗解精确求解.

对任意对称正定矩阵 $\bm M$, 记
\begin{align}
    \norm{\bm x}_{\bm M}&=\sqrt{\bm x^\T\bm M\bm x},
    \label{eq:notation-vector-energy}\\
    \norm{\bm X}_{\bm M,\mathrm{F}}&=\sqrt{\tr(\bm X^\T\bm M\bm X)},
    \label{eq:notation-left-frob}\\
    \avg{\bm X,\bm Y}_{\bm M,\mathrm{F}}&=\tr(\bm X^\T\bm M\bm Y),
    \label{eq:notation-left-inner}\\
    \norm{\bm X}_{\mathrm{F},\bm M}&=\sqrt{\tr(\bm X^\T\bm X\bm M)},
    \label{eq:notation-right-frob}\\
    \avg{\bm X,\bm Y}_{\mathrm{F},\bm M}&=\tr(\bm X^\T\bm Y\bm M),
    \label{eq:notation-right-inner}\\
    \norm{\bm X}_{\bm M}&=\norm{\bm M^{1/2}\bm X\bm M^{-1/2}}_2.
    \label{eq:notation-induced-energy}
\end{align}
此外, $\lambda_k(\cdot),\sigma_k(\cdot)$ 分别表示按非增序排列的第 $k$ 个特征值和奇异值, $\Ran(\cdot)$ 表示列空间, $\bm G^{\dagger_{\bm A}}=\bm A^{-1}\bm G^\T\bm A$. PCG 采用对角预条件共轭梯度法.

\section{能量极小插值与 PCG 迭代}
\label{sec:energy-pcg}

\subsection{能量极小值点}
\label{subsec:energy-endpoint}

定义插值能量与 Schur 补
\begin{align}
    J(\bm W)&=\frac{1}{2}\tr(\bm P(\bm W)^\T\bm A\bm P(\bm W)),
    \label{eq:energy-functional}\\
    \bm S&=\bm A_{\mathrm{CC}}-\bm C^\T\bm B^{-1}\bm C\succ \mathbf{O}.
    \label{eq:schur-complement}
\end{align}

\begin{proposition}
\label{prop:energy-minimizer}
插值能量的唯一极小值点为
\begin{align}
    \bm W^{\ast}&=-\bm B^{-1}\bm C.
    \label{eq:energy-minimizer}
\end{align}
令 $\bm\Delta=\bm W-\bm W^{\ast}$, 则
\begin{align}
    J(\bm W)-J(\bm W^{\ast})&=\frac{1}{2}\norm{\bm\Delta}_{\bm B,\mathrm{F}}^2,
    \label{eq:energy-gap}\\
    \bm P(\bm W)^\T\bm A\bm P(\bm W)&=\bm S+\bm\Delta^\T\bm B\bm\Delta.
    \label{eq:coarse-matrix-energy-error}
\end{align}
\end{proposition}

\begin{proof}
由\cref{eq:setup-block-A,eq:setup-interpolation} 作分块乘法,
\begin{align}
    \bm P(\bm W)^\T\bm A\bm P(\bm W)
    &=\bm W^\T\bm B\bm W+\bm W^\T\bm C+\bm C^\T\bm W+\bm A_{\mathrm{CC}},
    \label{eq:energy-block-matrix-expansion}
\end{align}
从而利用迹的性质,
\begin{align}
    J(\bm W)
    &=\frac12\tr(\bm W^\T\bm B\bm W)
      +\tr(\bm W^\T\bm C)
      +\frac12\tr(\bm A_{\mathrm{CC}}).
    \label{eq:energy-block-expansion}
\end{align}
对任意扰动 $\bm H$, 求 Fr\'echet 导数得到
\begin{align}
    J'(\bm W)[\bm H]
    &=\tr\left(\bm H^\T(\bm B\bm W+\bm C)\right).
    \label{eq:energy-frechet-derivative}
\end{align}
由\cref{eq:energy-frechet-derivative} 以及 $\bm B\succ\mathbf O$ 可知, $J$ 关于 $\bm W$ 严格凸, 因而唯一驻点 $\bm W^{\ast}$ 满足\cref{eq:energy-minimizer}.

由\cref{eq:energy-minimizer} 令 $\bm W=\bm W^\ast+\bm\Delta$. 展开能量差并使用 $\bm B\bm W^\ast+\bm C=\mathbf O$, 有
\begin{align}
    J(\bm W)-J(\bm W^\ast)
    &=\frac12\tr(\bm\Delta^\T\bm B\bm\Delta)
      +\tr\left(\bm\Delta^\T(\bm B\bm W^\ast+\bm C)\right)=\frac12\tr(\bm\Delta^\T\bm B\bm\Delta)
     =\frac12\norm{\bm\Delta}_{\bm B,\mathrm F}^2.
    \label{eq:energy-gap-proof-explicit}
\end{align}
同理, 将 $\bm W^\ast+\bm\Delta$ 代入\cref{eq:energy-block-matrix-expansion}, 并结合\cref{eq:schur-complement,eq:energy-minimizer}, 关于 $\bm\Delta$ 的线性项全部消掉, 而
\begin{align}
    {\bm W^\ast}^\T\bm B\bm W^\ast
    +{\bm W^\ast}^\T\bm C
    +\bm C^\T\bm W^\ast
    +\bm A_{\mathrm{CC}}
    &=\bm A_{\mathrm{CC}}-\bm C^\T\bm B^{-1}\bm C
     =\bm S.
    \label{eq:endpoint-coarse-matrix-proof}
\end{align}
故得到\cref{eq:coarse-matrix-energy-error}.
\end{proof}

\subsection{PCG 迭代}
\label{subsec:pcg-energy-geometry}

对\cref{eq:energy-minimizer} 的第 $j$ 列, 以几何插值为初值进行 PCG. 记
\begin{align}
    \bm B\bm w_j^{\ast}&=-\bm C\bm e_j,
    \label{eq:pcg-column-system}\\
    \bm e_{j,m}&=\bm w_j^{\ast}-\bm w_{j,m},\q
    \bm s_{j,m}=\bm w_{j,m+1}-\bm w_{j,m},
    \label{eq:pcg-column-error-step}\\
    \bm\Delta_m&=\bm W_m-\bm W^{\ast}.
    \label{eq:pcg-matrix-error}
\end{align}

\begin{proposition}
\label{prop:pcg-exact-identities}
对任意 PCG 步,
\begin{align}
    \norm{\bm\Delta_m}_{\bm B,\mathrm{F}}^2&=\norm{\bm\Delta_{m+1}}_{\bm B,\mathrm{F}}^2+\norm{\bm W_{m+1}-\bm W_m}_{\bm B,\mathrm{F}}^2.
    \label{eq:pcg-matrix-pythagoras}
\end{align}
若 $\bm W_{m+1}\ne\bm W_m$, 则
\begin{align}
    J(\bm W_{m+1})&<J(\bm W_m).
    \label{eq:pcg-energy-strict-decrease}
\end{align}
\end{proposition}

\begin{proof}
由\cref{eq:pcg-column-error-step}, $\bm e_{j,m}=\bm e_{j,m+1}+\bm s_{j,m}$. PCG 的 Galerkin 正交性给出 $\bm e_{j,m+1}\perp_{\bm B}\bm s_{j,m}$, 因此逐列有
\begin{align}
    \norm{\bm e_{j,m}}_{\bm B}^2
    &=\norm{\bm e_{j,m+1}}_{\bm B}^2
      +\norm{\bm s_{j,m}}_{\bm B}^2.
    \label{eq:pcg-column-pythagoras}
\end{align}
注意 $\bm\Delta_m$ 的第 $j$ 列为 $-\bm e_{j,m}$, 对 $j=1,\ldots,n_{\mathrm C}$ 求和即得\cref{eq:pcg-matrix-pythagoras}. 再结合\cref{eq:energy-gap}, 因而
\begin{align}
    J(\bm W_m)-J(\bm W_{m+1})
    &=\frac12\left(
        \norm{\bm\Delta_m}_{\bm B,\mathrm F}^2
        -\norm{\bm\Delta_{m+1}}_{\bm B,\mathrm F}^2
      \right)=\frac12\norm{\bm W_{m+1}-\bm W_m}_{\bm B,\mathrm F}^2.
    \label{eq:pcg-energy-decrease-explicit}
\end{align}
若该步非零, 右端严格为正, 从而得到\cref{eq:pcg-energy-strict-decrease}.
\end{proof}

\section{归一化误差与粗空间}
\label{sec:normalized-geometry}

令
\begin{align}
    \bm Z&=\bm B^{1/2}\bm\Delta\bm S^{-1/2},
    \label{eq:normalized-error}\\
    \bm Q_{\mathrm{F}}&=\begin{bNiceMatrix}
        \bm B^{-1/2}\\
        \mathbf{O}
    \end{bNiceMatrix},\q
    \bm Q_{\ast}=\bm P(\bm W^{\ast})\bm S^{-1/2}.
    \label{eq:A-orthonormal-basis}
\end{align}
由\cref{eq:setup-block-A,eq:A-orthonormal-basis,eq:energy-minimizer,eq:coarse-matrix-energy-error} 直接分块计算可得
\begin{align}
    \bm Q_{\mathrm F}^\T\bm A\bm Q_{\mathrm F}
    &=\bm B^{-1/2}\bm B\bm B^{-1/2}=\bm I,
    \label{eq:QF-A-QF}\\
    \bm Q_\ast^\T\bm A\bm Q_\ast
    &=\bm S^{-1/2}\bm P(\bm W^\ast)^\T\bm A\bm P(\bm W^\ast)\bm S^{-1/2}=\bm I,
    \label{eq:Qstar-A-Qstar}\\
    \bm Q_{\mathrm F}^\T\bm A\bm Q_\ast
    &=\bm B^{-1/2}(\bm B\bm W^\ast+\bm C)\bm S^{-1/2}=\mathbf O.
    \label{eq:QF-A-Qstar}
\end{align}
因此
\begin{align}
    \begin{bNiceMatrix}\bm Q_{\mathrm F}&\bm Q_\ast\end{bNiceMatrix}^\T
    \bm A
    \begin{bNiceMatrix}\bm Q_{\mathrm F}&\bm Q_\ast\end{bNiceMatrix}
    &=\bm I,
    \label{eq:A-orthonormal-combined}
\end{align}
即 $\Ran(\bm Q_{\mathrm F})\oplus\Ran(\bm Q_\ast)$ 构成一组 $\bm A$-单位正交基. 再令
\begin{align}
    \ti{\bm Q}_{\mathrm{F}}&=\bm A^{1/2}\bm Q_{\mathrm{F}},\q
    \ti{\bm Q}_{\ast}=\bm A^{1/2}\bm Q_{\ast},
    \label{eq:euclidean-orthonormal-basis}
\end{align}
则 $\begin{bNiceMatrix}\ti{\bm Q}_{\mathrm F}&\ti{\bm Q}_\ast\end{bNiceMatrix}$ 为欧氏正交矩阵.

\begin{proposition}
\label{prop:normalized-geometry}
对任意 $\bm W$, 有
\begin{align}
    J(\bm W)-J(\bm W^{\ast})&=\frac{1}{2}\norm{\bm Z}_{\mathrm{F},\bm S}^2,
    \label{eq:normalized-energy-gap}\\
    \bm P(\bm W)^\T\bm A\bm P(\bm W)&=\bm S^{1/2}(\bm I+\bm Z^\T\bm Z)\bm S^{1/2},
    \label{eq:normalized-coarse-matrix}\\
    \bm P(\bm W)&=\begin{bNiceMatrix}
        \bm Q_{\mathrm{F}} & \bm Q_{\ast}
    \end{bNiceMatrix}
    \begin{bNiceMatrix}
        \bm Z\\
        \bm I
    \end{bNiceMatrix}\bm S^{1/2}.
    \label{eq:interpolation-graph-factorization}
\end{align}
令 $\theta_k$ 为 $\Ran(\bm P(\bm W))$ 与能量极小粗空间 $\Ran(\bm P(\bm W^\ast))$ 关于 $\bm A$-内积的主角, 按非增序排列, 则 $\tan\theta_k=\sigma_k(\bm Z)$.
\end{proposition}

\begin{proof}
由\cref{eq:normalized-error},
\begin{align}
    \bm\Delta&=\bm B^{-1/2}\bm Z\bm S^{1/2},\q
    \bm\Delta^\T\bm B\bm\Delta
    =\bm S^{1/2}\bm Z^\T\bm Z\bm S^{1/2}.
    \label{eq:normalized-delta-substitution}
\end{align}
于是
\begin{align}
    \norm{\bm\Delta}_{\bm B,\mathrm F}^2
    &=\tr(\bm Z^\T\bm Z\bm S)
     =\norm{\bm Z}_{\mathrm F,\bm S}^2,
    \label{eq:normalized-norm-equivalence}
\end{align}
将\cref{eq:normalized-delta-substitution,eq:normalized-norm-equivalence} 分别代入\cref{eq:energy-gap,eq:coarse-matrix-energy-error}, 就得到\cref{eq:normalized-energy-gap,eq:normalized-coarse-matrix}.

另一方面, 由\cref{eq:A-orthonormal-basis,eq:normalized-delta-substitution},
\begin{align}
    \bm P(\bm W)=\bm P(\bm W^\ast)
      +\begin{bNiceMatrix}\bm\Delta\\\mathbf O\end{bNiceMatrix}=\bm Q_\ast\bm S^{1/2}
      +\bm Q_{\mathrm F}\bm Z\bm S^{1/2}=\begin{bNiceMatrix}\bm Q_{\mathrm F}&\bm Q_\ast\end{bNiceMatrix}
      \begin{bNiceMatrix}\bm Z\\\bm I\end{bNiceMatrix}\bm S^{1/2},
    \label{eq:interpolation-graph-factorization-proof}
\end{align}
故在 $\bm A$-单位正交坐标中, 当前粗空间恰是图空间
\begin{align}
    \mathcal G(\bm Z)&:=\Ran\begin{bNiceMatrix}\bm Z\\\bm I\end{bNiceMatrix},
    \label{eq:graph-space-definition}
\end{align}
而能量极小粗空间对应 $\mathcal G(\mathbf O)$.

再取 $\bm Z^\T\bm Z$ 的单位正交特征向量 $\bm v_k$, 满足 $\bm Z^\T\bm Z\bm v_k=\sigma_k(\bm Z)^2\bm v_k$. 图空间中对应的单位向量可取为
\begin{align}
    \bm q_k
    &=\frac{1}{\sqrt{1+\sigma_k(\bm Z)^2}}
      \begin{bNiceMatrix}\bm Z\bm v_k\\\bm v_k\end{bNiceMatrix},
    \label{eq:graph-principal-vector}
\end{align}
而 $\mathcal G(\mathbf O)$ 中对应的单位向量为
\begin{align}
    \bm q_k^\ast&=\begin{bNiceMatrix}\mathbf O\\\bm v_k\end{bNiceMatrix}.
    \label{eq:endpoint-principal-vector}
\end{align}
由\cref{eq:graph-principal-vector,eq:endpoint-principal-vector}, $\{\bm q_k\}$ 与 $\{\bm q_k^\ast\}$ 分别构成两个子空间的单位正交基, 且交叉 Gram 矩阵在该配对下为对角阵, 第 $k$ 个对角元为 $(1+\sigma_k(\bm Z)^2)^{-1/2}$. 因而相应主角满足 $\tan\theta_k=\sigma_k(\bm Z)$.
\end{proof}

\section{两网格收敛率}
\label{sec:tg-objective}

\subsection{能量坐标与 Gauss--Seidel 迭代}
\label{subsec:energy-coordinate-tg}

定义前磨光子在 $\bm A$-能量坐标下的表示与粗空间
\begin{align}
    \bm T&=\bm A^{1/2}\bm G\bm A^{-1/2},
    \label{eq:smoother-energy-coordinate}\\
    \mathcal U&=\Ran(\bm A^{1/2}\bm P(\bm W)).
    \label{eq:energy-coordinate-coarse-space}
\end{align}
到 $\Ran(\bm P(\bm W))$ 的 $\bm A$-正交投影记为
\begin{align}
    \bm Q_{\mathcal U}^{\bm A}
    &=\bm P(\bm W)
      (\bm P(\bm W)^\T\bm A\bm P(\bm W))^{-1}
      \bm P(\bm W)^\T\bm A.
    \label{eq:A-projector}
\end{align}
在能量坐标下,
\begin{align}
    \bm\Pi_{\mathcal U}:=\bm A^{1/2}\bm Q_{\mathcal U}^{\bm A}\bm A^{-1/2}=\bm A^{1/2}\bm P(\bm W)
      (\bm P(\bm W)^\T\bm A\bm P(\bm W))^{-1}
      \bm P(\bm W)^\T\bm A^{1/2},
    \label{eq:energy-coordinate-projector}
\end{align}
由\cref{eq:A-projector,eq:energy-coordinate-projector}, $\bm\Pi_{\mathcal U}$ 是到 $\mathcal U$ 的欧氏正交投影.

两网格误差传播算子为
\begin{align}
    \bm E_{\mathrm{TG}}
    &=\bm G^{\dagger_{\bm A}}
      (\bm I-\bm Q_{\mathcal U}^{\bm A})\bm G.
    \label{eq:tg-error-propagator}
\end{align}
由\cref{eq:smoother-energy-coordinate} 和 $\bm G^{\dagger_{\bm A}}=\bm A^{-1}\bm G^\T\bm A$,
\begin{align}
    \bm A^{1/2}\bm G^{\dagger_{\bm A}}\bm A^{-1/2}
    &=\bm A^{-1/2}\bm G^\T\bm A^{1/2}
     =\bm T^\T,
    \label{eq:adjoint-smoother-energy-coordinate}
\end{align}
并且由\cref{eq:energy-coordinate-projector},
\begin{align}
    \bm A^{1/2}(\bm I-\bm Q_{\mathcal U}^{\bm A})\bm A^{-1/2}
    &=\bm I-\bm\Pi_{\mathcal U}.
    \label{eq:coarse-correction-energy-coordinate}
\end{align}
结合\cref{eq:tg-error-propagator,eq:adjoint-smoother-energy-coordinate,eq:coarse-correction-energy-coordinate},
\begin{align}
    \bm A^{1/2}\bm E_{\mathrm{TG}}\bm A^{-1/2}
    &=\bm T^\T(\bm I-\bm\Pi_{\mathcal U})\bm T.
    \label{eq:tg-energy-coordinate}
\end{align}
由于 $\bm I-\bm\Pi_{\mathcal U}$ 是对称幂等矩阵,
\begin{align}
    \bm T^\T(\bm I-\bm\Pi_{\mathcal U})\bm T
    &=\bigl((\bm I-\bm\Pi_{\mathcal U})\bm T\bigr)^\T
      \bigl((\bm I-\bm\Pi_{\mathcal U})\bm T\bigr)\succeq\mathbf O.
    \label{eq:tg-energy-coordinate-gram}
\end{align}
由\cref{eq:tg-energy-coordinate,eq:tg-energy-coordinate-gram}, 得到
\begin{align}
    \rho(\bm E_{\mathrm{TG}})
    &=\norm{(\bm I-\bm\Pi_{\mathcal U})\bm T}_2^2
     =:\rho_{\mathrm{TG}}(\mathcal U).
    \label{eq:tg-objective-exact}
\end{align}

\begin{lemma}
\label{lem:gs-nonexpansive}
前向 Gauss--Seidel 迭代满足
\begin{align}
    \norm{\bm T}_2&=\norm{\bm G}_{\bm A}\le 1.
    \label{eq:gs-nonexpansive}
\end{align}
\end{lemma}

\begin{proof}
一次坐标修正为
\begin{align}
    \bm e'&=\bm e-\frac{\bm\eps_i^\T\bm A\bm e}{a_{ii}}\bm\eps_i,
    \label{eq:gs-coordinate-update}
\end{align}
且
\begin{align}
    \norm{\bm e'}_{\bm A}^2&=\norm{\bm e}_{\bm A}^2-\frac{(\bm\eps_i^\T\bm A\bm e)^2}{a_{ii}}\le\norm{\bm e}_{\bm A}^2.
    \label{eq:gs-coordinate-energy}
\end{align}
一次前向 Gauss--Seidel 迭代是\cref{eq:gs-coordinate-update} 的连续复合, 故 $\norm{\bm G\bm e}_{\bm A}\le\norm{\bm e}_{\bm A}$; 结合\cref{eq:notation-induced-energy,eq:smoother-energy-coordinate} 即得\cref{eq:gs-nonexpansive}.
\end{proof}

\subsection{图空间表示}
\label{subsec:graph-projector-stability}

在\cref{eq:euclidean-orthonormal-basis} 的正交基底下, 定义
\begin{align}
    \begin{bNiceMatrix}
        \bm T_{\mathrm F}\\
        \bm T_\ast
    \end{bNiceMatrix}
    &=\begin{bNiceMatrix}
        \ti{\bm Q}_{\mathrm F}^\T\\
        \ti{\bm Q}_{\ast}^\T
    \end{bNiceMatrix}\bm T.
    \label{eq:T-block-coordinates}
\end{align}
由\cref{eq:gs-nonexpansive} 有 $\norm{\bm T_{\mathrm F}}_2,\norm{\bm T_\ast}_2\le1$.

\begin{proposition}
\label{prop:graph-projector}
在\cref{eq:euclidean-orthonormal-basis} 的坐标下, 图空间 $\Ran\begin{bNiceMatrix}\bm Z\\\bm I\end{bNiceMatrix}$ 的正交补具有欧氏单位正交基
\begin{align}
    \bm N(\bm Z)&=\begin{bNiceMatrix}
        \bm I\\
        -\bm Z^\T
    \end{bNiceMatrix}(\bm I+\bm Z\bm Z^\T)^{-1/2}.
    \label{eq:graph-complement-basis}
\end{align}
因此
\begin{align}
    \rho_{\mathrm{TG}}(\bm Z)&=\norm*{(\bm I+\bm Z\bm Z^\T)^{-1/2}
    (\bm T_{\mathrm F}-\bm Z\bm T_\ast)}_2^2.
    \label{eq:tg-objective-graph}
\end{align}
\end{proposition}

\begin{proof}
记 $\ti{\bm Q}=\begin{bNiceMatrix}\ti{\bm Q}_{\mathrm F}&\ti{\bm Q}_\ast\end{bNiceMatrix}$. 由\cref{eq:interpolation-graph-factorization-proof}, 在该正交坐标中当前粗空间就是 $\mathcal G(\bm Z)$. 先直接验证
\begin{align}
    \begin{bNiceMatrix}\bm Z^\T&\bm I\end{bNiceMatrix}\bm N(\bm Z)
    &=\bigl(\bm Z^\T-\bm Z^\T\bigr)
      (\bm I+\bm Z\bm Z^\T)^{-1/2}
     =\mathbf O,
    \label{eq:graph-complement-orthogonality}\\
    \bm N(\bm Z)^\T\bm N(\bm Z)
    &=(\bm I+\bm Z\bm Z^\T)^{-1/2}
      (\bm I+\bm Z\bm Z^\T)
      (\bm I+\bm Z\bm Z^\T)^{-1/2}
     =\bm I.
    \label{eq:graph-complement-normalization}
\end{align}
故 $\bm N(\bm Z)$ 的列构成 $\mathcal G(\bm Z)^\perp$ 的单位正交基, 从而
\begin{align}
    \bm I-\bm\Pi_{\mathcal U}
    &=\ti{\bm Q}\bm N(\bm Z)\bm N(\bm Z)^\T\ti{\bm Q}^\T.
    \label{eq:graph-complement-projector}
\end{align}
利用\cref{eq:tg-objective-exact,eq:T-block-coordinates,eq:graph-complement-projector} 以及 $\ti{\bm Q}$ 的正交性,
\begin{align}
    \rho_{\mathrm{TG}}(\bm Z)
    &=\norm{\bm N(\bm Z)^\T\ti{\bm Q}^\T\bm T}_2^2=\norm*{
        (\bm I+\bm Z\bm Z^\T)^{-1/2}
        \begin{bNiceMatrix}\bm I&-\bm Z\end{bNiceMatrix}
        \begin{bNiceMatrix}\bm T_{\mathrm F}\\\bm T_\ast\end{bNiceMatrix}
      }_2^2=\norm*{
        (\bm I+\bm Z\bm Z^\T)^{-1/2}
        (\bm T_{\mathrm F}-\bm Z\bm T_\ast)
      }_2^2,
    \label{eq:tg-objective-coordinate-projection}
\end{align}
即得到\cref{eq:tg-objective-graph}.
\end{proof}

\section{能量极小值点附近的局部谱分析}
\label{sec:local-spectral-analysis}

\subsection{最大奇异值的二阶扰动界}
\label{subsec:singular-value-perturbation}

\begin{lemma}
\label{lem:singular-value-perturbation}
设
\begin{align}
    \sigma&=\sigma_1(\bm M)>\sigma_2(\bm M),\q
    \dt=\sigma_1(\bm M)-\sigma_2(\bm M)>0,
    \label{eq:sv-gap}\\
    \bm M\bm v&=\sigma\bm u,\q \bm M^\T\bm u=\sigma\bm v,
    \label{eq:sv-singular-vectors}
\end{align}
其中 $\bm u,\bm v$ 为单位奇异向量. 若 $\norm{\bm E}_2\le\dt/4$, 则
\begin{align}
    \abs*{\sigma_1(\bm M+\bm E)-\sigma-\bm u^\T\bm E\bm v}
    &\le\frac{2}{\dt}\norm{\bm E}_2^2.
    \label{eq:sv-second-order-bound}
\end{align}
\end{lemma}

\begin{proof}
作对称扩张
\begin{align}
    \bm M_0&=
    \begin{bNiceMatrix}\mathbf O&\bm M\\\bm M^\T&\mathbf O\end{bNiceMatrix},\q
    \bm E_0=
    \begin{bNiceMatrix}\mathbf O&\bm E\\\bm E^\T&\mathbf O\end{bNiceMatrix},\q
    \bm q=\frac1{\sqrt2}
    \begin{bNiceMatrix}\bm u\\\bm v\end{bNiceMatrix}.
    \label{eq:sv-symmetric-dilation}
\end{align}
由\cref{eq:sv-gap,eq:sv-singular-vectors,eq:sv-symmetric-dilation}, 于是 $\bm M_0\bm q=\sigma\bm q$, $\norm{\bm E_0}_2=\norm{\bm E}_2$, 且
\begin{align}
    \lambda_1(\bm M_0+\bm E_0)
    &=\sigma_1(\bm M+\bm E)=:\sigma'.
    \label{eq:sv-dilation-top-eigenvalue}
\end{align}
取 $\sigma'$ 的单位特征向量并作正交分解
\begin{align}
    \bm q'&=\gamma\bm q+\bm z_\perp,\q
    \bm z_\perp\perp\bm q.
    \label{eq:sv-eigenvector-decomposition}
\end{align}
Weyl 不等式给出 $\sigma'\ge\sigma-\norm{\bm E}_2$. 记 $\bm Q_\perp=\bm I-\bm q\bm q^\T$. 将
\begin{align}
    (\bm M_0+\bm E_0)\bm q'&=\sigma'\bm q'
    \label{eq:sv-perturbed-eigen-equation}
\end{align}
投影到 $\bm q^\perp$. 由于 $\bm M_0$ 对称且 $\bm q$ 是其特征向量, $\bm q^\perp$ 在 $\bm M_0$ 下不变, 因而
\begin{align}
    (\sigma'\bm I-\bm M_0)\bm z_\perp
    &=\bm Q_\perp\bm E_0(\gamma\bm q+\bm z_\perp).
    \label{eq:sv-projected-eigen-equation}
\end{align}
结合\cref{eq:sv-gap,eq:sv-projected-eigen-equation}, $\bm M_0$ 在 $\bm q^\perp$ 上的最大特征值不超过 $\sigma_2(\bm M)$, 所以
\begin{align}
    (\dt-\norm{\bm E}_2)\norm{\bm z_\perp}_2\le(\sigma'-\sigma_2(\bm M))\norm{\bm z_\perp}_2\le\norm{(\sigma'\bm I-\bm M_0)\bm z_\perp}_2\le\norm{\bm E}_2
       \bigl(\abs\gamma+\norm{\bm z_\perp}_2\bigr).
    \label{eq:sv-eigenvector-prebound}
\end{align}
若 $\gamma=0$, 则 $\norm{\bm z_\perp}_2=1$, 上式推出 $\dt\le2\norm{\bm E}_2$, 与 $\norm{\bm E}_2\le\dt/4$ 矛盾. 因而 $\gamma\ne0$. 将含 $\norm{\bm z_\perp}_2$ 的项移到左侧,
\begin{align}
    \frac{\norm{\bm z_\perp}_2}{\abs\gamma}
    &\le\frac{\norm{\bm E}_2}{\dt-2\norm{\bm E}_2}.
    \label{eq:sv-eigenvector-tilt}
\end{align}

再将特征方程左乘 $\bm q^\T$. 利用 $\bm q^\T\bm M_0=\sigma\bm q^\T$ 与 $\bm q^\T\bm z_\perp=0$,
\begin{align}
    (\sigma'-\sigma)\gamma
    &=\gamma\bm q^\T\bm E_0\bm q
      +\bm q^\T\bm E_0\bm z_\perp.
    \label{eq:sv-top-eigenvalue-projection}
\end{align}
而
\begin{align}
    \bm q^\T\bm E_0\bm q
    &=\bm u^\T\bm E\bm v.
    \label{eq:sv-dilation-first-order-term}
\end{align}
结合\cref{eq:sv-eigenvector-tilt,eq:sv-top-eigenvalue-projection,eq:sv-dilation-first-order-term}, 于是
\begin{align}
    \abs*{\sigma'-\sigma-\bm u^\T\bm E\bm v}=\frac{\abs{\bm q^\T\bm E_0\bm z_\perp}}{\abs\gamma}\le\norm{\bm E}_2
      \frac{\norm{\bm z_\perp}_2}{\abs\gamma}\le\frac{\norm{\bm E}_2^2}{\dt-2\norm{\bm E}_2}\le\frac{2}{\dt}\norm{\bm E}_2^2,
    \label{eq:sv-second-order-bound-proof}
\end{align}
即证.
\end{proof}

\subsection{能量极小值点处的局部展开}
\label{subsec:endpoint-local-expansion}

记 $f(\bm Z)=\sqrt{\rho_{\mathrm{TG}}(\bm Z)}$. 由\cref{eq:tg-objective-graph},
\begin{align}
    f(\bm Z)&=\norm{\bm L(\bm Z)}_2,\q
    \bm L(\bm Z):=(\bm I+\bm Z\bm Z^\T)^{-1/2}(\bm T_{\mathrm F}-\bm Z\bm T_\ast).
    \label{eq:local-L-definition}
\end{align}
假设端点最大奇异值简单,
\begin{align}
    \sigma&:=\sigma_1(\bm T_{\mathrm F})>\sigma_2(\bm T_{\mathrm F}),\q
    \dt:=\sigma_1(\bm T_{\mathrm F})-\sigma_2(\bm T_{\mathrm F})>0,
    \label{eq:endpoint-singular-gap}
\end{align}
并取对应单位左右奇异向量满足
\begin{align}
    \bm T_{\mathrm F}\bm v&=\sigma\bm u_{\mathrm F},\q
    \bm T_{\mathrm F}^\T\bm u_{\mathrm F}=\sigma\bm v,\q
    \bm y:=\bm T_\ast\bm v.
    \label{eq:endpoint-singular-vectors}
\end{align}
由 $\norm{\bm T_{\mathrm F}}_2\le1$ 还知 $0<\dt\le\sigma\le1$.

\begin{theorem}
\label{thm:endpoint-local-expansion}
若
\begin{align}
    \norm{\bm Z}_{\mathrm F}&\le \dt/8,
    \label{eq:local-ball-condition}
\end{align}
则
\begin{align}
    f(\bm Z)&=\sigma-\bm u_{\mathrm F}^\T\bm Z\bm y+R(\bm Z),
    \label{eq:local-expansion}\\
    \abs{R(\bm Z)}&\le\left(1+8/\dt\right)\norm{\bm Z}_{\mathrm F}^2.
    \label{eq:local-remainder-bound}
\end{align}
\end{theorem}

\begin{proof}
令
\begin{align}
    \bm A(\bm Z)&:=(\bm I+\bm Z\bm Z^\T)^{-1/2},\q
    z:=\norm{\bm Z}_2.
    \label{eq:local-normalization-matrix}
\end{align}
因为 $\bm Z\bm Z^\T\succeq\mathbf O$,
\begin{align}
    \norm{\bm A(\bm Z)-\bm I}_2
    &=1-\frac1{\sqrt{1+z^2}}
     \le\frac12z^2.
    \label{eq:local-normalization-bound}
\end{align}
由\cref{eq:local-L-definition} 将 $\bm L(\bm Z)$ 分成端点矩阵与扰动:
\begin{align}
    \bm L(\bm Z)
    &=\bm T_{\mathrm F}+\bm E(\bm Z),
    \label{eq:local-L-perturbation-split}\\
    \bm E(\bm Z)
    &=-\bm Z\bm T_\ast+\bm D(\bm Z),
    \label{eq:local-E-definition}\\
    \bm D(\bm Z)
    &:=(\bm A(\bm Z)-\bm I)
      (\bm T_{\mathrm F}-\bm Z\bm T_\ast).
    \label{eq:local-D-definition}
\end{align}
由\cref{eq:gs-nonexpansive,eq:T-block-coordinates,eq:local-ball-condition,eq:local-normalization-bound} 以及 $z\le\norm{\bm Z}_{\mathrm F}\le\dt/8\le1$,
\begin{align}
    \norm{\bm D(\bm Z)}_2
    &\le\frac12z^2(1+z)
     \le z^2,
    \label{eq:local-D-bound}\\
    \norm{\bm E(\bm Z)}_2
    &\le z+z^2
     \le2z
     \le \dt/4.
    \label{eq:local-E-bounds}
\end{align}
因此可对 $\bm T_{\mathrm F}+\bm E(\bm Z)$ 使用\cref{lem:singular-value-perturbation}:
\begin{align}
    \abs*{
        f(\bm Z)-\sigma
        -\bm u_{\mathrm F}^\T\bm E(\bm Z)\bm v
    }
    &\le\frac2\dt\norm{\bm E(\bm Z)}_2^2.
    \label{eq:local-apply-sv-perturbation}
\end{align}
注意
\begin{align}
    \bm u_{\mathrm F}^\T\bm E(\bm Z)\bm v=-\bm u_{\mathrm F}^\T\bm Z\bm T_\ast\bm v
      +\bm u_{\mathrm F}^\T\bm D(\bm Z)\bm v=-\bm u_{\mathrm F}^\T\bm Z\bm y
      +\bm u_{\mathrm F}^\T\bm D(\bm Z)\bm v.
    \label{eq:local-first-order-separation}
\end{align}
于是
\begin{align}
    \abs*{f(\bm Z)-\sigma+\bm u_{\mathrm F}^\T\bm Z\bm y}\le\abs{\bm u_{\mathrm F}^\T\bm D(\bm Z)\bm v}
      +\frac2\dt\norm{\bm E(\bm Z)}_2^2\le z^2+\frac8\dt z^2\le\left(1+\frac8\dt\right)\norm{\bm Z}_{\mathrm F}^2.
    \label{eq:local-expansion-bound-proof}
\end{align}
定义
\begin{align}
    R(\bm Z)&:=f(\bm Z)-\sigma+\bm u_{\mathrm F}^\T\bm Z\bm y,
    \label{eq:local-remainder-definition}
\end{align}
即得到\cref{eq:local-expansion,eq:local-remainder-bound}.
\end{proof}

\section{PCG 迭代的相邻谱半径变化}
\label{sec:pcg-step-performance}

令
\begin{align}
    \bm D_m&:=\bm Z_m-\bm Z_{m+1},\q
    r_m:=\norm{\bm Z_m}_{\mathrm F,\bm S}.
    \label{eq:path-step-and-radius}
\end{align}
由\cref{eq:normalized-error,eq:pcg-matrix-pythagoras,eq:normalized-energy-gap} 可得
\begin{align}
    r_m^2&=r_{m+1}^2+\norm{\bm D_m}_{\mathrm F,\bm S}^2,
    \label{eq:path-radius-pythagoras}\\
    J(\bm W_m)-J(\bm W_{m+1})&=\frac12\norm{\bm D_m}_{\mathrm F,\bm S}^2.
    \label{eq:path-energy-decrease}
\end{align}

\begin{lemma}
\label{lem:pcg-one-step-contraction}
令
\begin{align}
    \bm D_{\bm B}&=\diag(\bm B),\q
    \kappa=\kappa_2(\bm D_{\bm B}^{-1/2}\bm B\bm D_{\bm B}^{-1/2}),
    \label{eq:jacobi-condition-number}\\
    c_1&=\frac{\kappa-1}{\kappa+1},\q
    c_2=\frac{2\sqrt\kappa}{\kappa+1}.
    \label{eq:pcg-c1-c2}
\end{align}
则
\begin{align}
    r_{m+1}&\le c_1r_m,
    \label{eq:pcg-radius-contraction}\\
    \norm{\bm D_m}_{\mathrm F,\bm S}&\ge c_2r_m.
    \label{eq:pcg-step-lower-bound}
\end{align}
\end{lemma}

\begin{proof}
PCG 在相同 Krylov 子空间中最小化 $\bm B$-能量误差, 因而一步误差不大于以同一对角预条件器作一步最速下降. 预条件最速下降的标准收缩界逐列给出
\begin{align}
    \norm{\bm e_{j,m+1}}_{\bm B}
    &\le\frac{\kappa-1}{\kappa+1}
      \norm{\bm e_{j,m}}_{\bm B}
     =c_1\norm{\bm e_{j,m}}_{\bm B}.
    \label{eq:pcg-column-one-step-contraction}
\end{align}
对所有列平方求和, 并利用 $r_m=\norm{\bm\Delta_m}_{\bm B,\mathrm F}$, 得
\begin{align}
    r_{m+1}^2
    &\le c_1^2r_m^2.
    \label{eq:pcg-radius-contraction-squared}
\end{align}
因此得到\cref{eq:pcg-radius-contraction}. 再由\cref{eq:path-radius-pythagoras},
\begin{align}
    \norm{\bm D_m}_{\mathrm F,\bm S}^2=r_m^2-r_{m+1}^2\ge(1-c_1^2)r_m^2
     =\frac{4\kappa}{(\kappa+1)^2}r_m^2
     =c_2^2r_m^2,
    \label{eq:pcg-step-lower-bound-proof}
\end{align}
即得\cref{eq:pcg-step-lower-bound}.
\end{proof}

若 $\bm D_m\ne\mathbf O$, 定义一步对齐量与判据常数
\begin{align}
    \nu_m&:=\frac{\bm u_{\mathrm F}^\T\bm D_m\bm y}{\norm{\bm D_m}_{\mathrm F,\bm S}},
    \label{eq:alignment-nu}\\
    C&:=\frac{(1+8/\dt)(1+c_1^2)}{\lambda_{\min}(\bm S)c_2}.
    \label{eq:step-criterion-C}
\end{align}

\begin{theorem}
\label{thm:pcg-two-sided-step-criterion}
设 $\bm D_m\ne\mathbf O$, 且 $\bm Z_m,\bm Z_{m+1}$ 均满足\cref{eq:local-ball-condition}. 则
\begin{align}
    J(\bm W_{m+1})&<J(\bm W_m),
    \label{eq:step-energy-strict-decrease}\\
    \nu_m>Cr_m&\q\Longrightarrow\q
    \rho_{\mathrm{TG}}(\bm Z_{m+1})>\rho_{\mathrm{TG}}(\bm Z_m),
    \label{eq:step-deterioration-condition}\\
    \nu_m<-Cr_m&\q\Longrightarrow\q
    \rho_{\mathrm{TG}}(\bm Z_{m+1})<\rho_{\mathrm{TG}}(\bm Z_m).
    \label{eq:step-improvement-condition}
\end{align}
\end{theorem}

\begin{proof}
先由\cref{eq:path-energy-decrease},
\begin{align}
    J(\bm W_m)-J(\bm W_{m+1})
    &=\frac12\norm{\bm D_m}_{\mathrm F,\bm S}^2>0,
    \label{eq:step-energy-strict-proof}
\end{align}
故能量严格下降.

记 $K:=1+8/\dt$. 由\cref{eq:local-expansion,eq:local-remainder-bound} 的局部展开,
\begin{align}
    f(\bm Z_{m+1})-f(\bm Z_m)
    &=\bm u_{\mathrm F}^\T(\bm Z_m-\bm Z_{m+1})\bm y
      +R(\bm Z_{m+1})-R(\bm Z_m)\\
    &=\bm u_{\mathrm F}^\T\bm D_m\bm y
      +R(\bm Z_{m+1})-R(\bm Z_m).
    \label{eq:step-f-difference-explicit}
\end{align}
余项满足
\begin{align}
    \abs{R(\bm Z_{m+1})-R(\bm Z_m)}
    &\le K\left(
        \norm{\bm Z_m}_{\mathrm F}^2
        +\norm{\bm Z_{m+1}}_{\mathrm F}^2
      \right).
    \label{eq:step-linearization-error}
\end{align}
又因为
\begin{align}
    r_m^2
    &=\tr(\bm Z_m^\T\bm Z_m\bm S)
     \ge\lambda_{\min}(\bm S)\norm{\bm Z_m}_{\mathrm F}^2,
    \label{eq:step-Zm-by-rm}
\end{align}
以及\cref{eq:pcg-radius-contraction}, 所以
\begin{align}
    \norm{\bm Z_m}_{\mathrm F}^2
    +\norm{\bm Z_{m+1}}_{\mathrm F}^2
    &\le\frac{1+c_1^2}{\lambda_{\min}(\bm S)}r_m^2.
    \label{eq:step-Z-bound-by-r}
\end{align}

若 $\nu_m>Cr_m$, 由\cref{eq:alignment-nu,eq:pcg-step-lower-bound,eq:step-criterion-C},
\begin{align}
    f(\bm Z_{m+1})-f(\bm Z_m)
    &\ge\nu_m\norm{\bm D_m}_{\mathrm F,\bm S}
      -\frac{K(1+c_1^2)}{\lambda_{\min}(\bm S)}r_m^2\\
    &\ge c_2\nu_m r_m
      -\frac{K(1+c_1^2)}{\lambda_{\min}(\bm S)}r_m^2\\
    &=c_2r_m(\nu_m-Cr_m)>0.
    \label{eq:step-deterioration-proof-chain}
\end{align}
由于 $f=\sqrt{\rho_{\mathrm{TG}}}\ge0$, 这等价于 $\rho_{\mathrm{TG}}(\bm Z_{m+1})>\rho_{\mathrm{TG}}(\bm Z_m)$.

若 $\nu_m<-Cr_m$, 则 $\nu_m<0$. 此时
\begin{align}
    \nu_m\norm{\bm D_m}_{\mathrm F,\bm S}
    &\le c_2\nu_m r_m.
    \label{eq:negative-alignment-step-bound}
\end{align}
于是
\begin{align}
    f(\bm Z_{m+1})-f(\bm Z_m)
    &\le\nu_m\norm{\bm D_m}_{\mathrm F,\bm S}
      +\frac{K(1+c_1^2)}{\lambda_{\min}(\bm S)}r_m^2\\
    &\le c_2\nu_m r_m
      +\frac{K(1+c_1^2)}{\lambda_{\min}(\bm S)}r_m^2\\
    &=c_2r_m(\nu_m+Cr_m)<0,
    \label{eq:step-improvement-proof-chain}
\end{align}
故 $\rho_{\mathrm{TG}}(\bm Z_{m+1})<\rho_{\mathrm{TG}}(\bm Z_m)$. 即证\cref{eq:step-deterioration-condition,eq:step-improvement-condition}.
\end{proof}

由\cref{eq:pcg-radius-contraction} 可知 $Cr_m$ 沿非终止 PCG 迭代指数衰减, 因而\cref{eq:step-deterioration-condition,eq:step-improvement-condition} 中无法判定符号的区间随 $m$ 收缩.

\subsection{判据常数的量级估计}
\label{subsec:criterion-scaling}

以下考虑固定有界区域上 Dirichlet 扩散问题的标准二维五点差分离散, 并假设扩散系数满足
\begin{align}
    0<\alpha&\le a(x)\le\beta,\q
    \chi:=\frac{\beta}{\alpha}\ge1.
    \label{eq:coefficient-bounds-and-contrast}
\end{align}
由\cref{eq:pcg-c1-c2} 可得
\begin{align}
    \frac{1+c_1^2}{c_2}
    &=\frac{\kappa^2+1}{\sqrt\kappa(\kappa+1)}
     =\Theta(\sqrt\kappa).
    \label{eq:pcg-factor-scaling}
\end{align}

对任意 $\bm x_{\mathrm C}$, 有
\begin{align}
    \bm x_{\mathrm C}^\T\bm S\bm x_{\mathrm C}
    &=\min_{\bm x_{\mathrm F}}
      \begin{bNiceMatrix}\bm x_{\mathrm F}^\T&\bm x_{\mathrm C}^\T\end{bNiceMatrix}
      \bm A
      \begin{bNiceMatrix}\bm x_{\mathrm F}\\\bm x_{\mathrm C}\end{bNiceMatrix}.
    \label{eq:schur-variational-characterization}
\end{align}
因此
\begin{align}
    \bm x_{\mathrm C}^\T\bm S\bm x_{\mathrm C}\ge\lambda_{\min}(\bm A)
      \min_{\bm x_{\mathrm F}}
      \left(\norm{\bm x_{\mathrm F}}_2^2+\norm{\bm x_{\mathrm C}}_2^2\right)=\lambda_{\min}(\bm A)\norm{\bm x_{\mathrm C}}_2^2.
    \label{eq:schur-lower-derivation}
\end{align}
对固定有界区域上的标准 Dirichlet 五点扩散离散, $\lambda_{\min}(\bm A)\gtrsim\alpha$, 故
\begin{align}
    \lambda_{\min}(\bm S)&\gtrsim\alpha.
    \label{eq:schur-lambda-lower}
\end{align}

再估计对角预条件后的条件数. 五点扩散矩阵的细点主块满足
\begin{align}
    \bm B&\preceq2\bm D_{\bm B},\q
    \lambda_{\max}(\bm D_{\bm B})\lesssim\beta h^{-2}.
    \label{eq:B-diagonal-comparison}
\end{align}
前者给出
\begin{align}
    \lambda_{\max}(\bm D_{\bm B}^{-1/2}\bm B\bm D_{\bm B}^{-1/2})
    &\le2.
    \label{eq:jacobi-preconditioned-lambda-max}
\end{align}
另一方面, $\bm B$ 是 $\bm A$ 的主子矩阵, 故
$\lambda_{\min}(\bm B)\ge\lambda_{\min}(\bm A)\gtrsim\alpha$. 对任意非零 $\bm x$,
\begin{align}
    \frac{\bm x^\T\bm B\bm x}{\bm x^\T\bm D_{\bm B}\bm x}
    &\ge\frac{\lambda_{\min}(\bm B)}{\lambda_{\max}(\bm D_{\bm B})}
     \gtrsim\frac{\alpha}{\beta h^{-2}}
     =\chi^{-1}h^2.
    \label{eq:jacobi-preconditioned-lambda-min}
\end{align}
结合\cref{eq:jacobi-preconditioned-lambda-max,eq:jacobi-preconditioned-lambda-min}, 因此
\begin{align}
    \kappa
    &=\kappa_2(\bm D_{\bm B}^{-1/2}\bm B\bm D_{\bm B}^{-1/2})
     =\O(\chi h^{-2}).
    \label{eq:kappa-scaling}
\end{align}
最后将以上估计代入判据常数 $C$,
\begin{align}
    C
    &=\frac{1+8/\dt}{\lambda_{\min}(\bm S)}
      \frac{1+c_1^2}{c_2}
     =\O\left(\frac{\sqrt\chi}{\alpha h\dt}\right).
    \label{eq:C-scaling}
\end{align}
该估计说明判据常数同时受网格尺度、对比度和端点谱间隙影响.

\section{局部判据的渐近可达性}
\label{sec:asymptotic-entry}

本节继续采用\cref{subsec:criterion-scaling}的标准二维五点差分设置. 对固定有限维问题, 以下 $m$ 均理解为 PCG 精确终止前的非平凡迭代步; 若 PCG 更早精确终止, 则已直接到达 $\bm Z=\mathbf O$.

在\cref{eq:jacobi-condition-number} 的同一预条件系统上, 标准 PCG 估计给出
\begin{align}
    c_3&:=\frac{\sqrt\kappa-1}{\sqrt\kappa+1},\q
    r_m\le2c_3^m r_0.
    \label{eq:pcg-global-matrix-bound}
\end{align}
再结合\cref{eq:path-step-and-radius,eq:notation-right-frob,eq:schur-lambda-lower,eq:pcg-global-matrix-bound},
\begin{align}
    \norm{\bm Z_m}_{\mathrm F}
    &\le\frac{2r_0}{\sqrt{\lambda_{\min}(\bm S)}}c_3^m.
    \label{eq:Z-global-exponential-bound}
\end{align}

\begin{assumption}
\label{ass:polynomial-data}
沿所考虑的离散问题族, 存在固定 $p\ge0$ 使得
\begin{align}
    \dt^{-1}=\O\left((\chi h^{-1})^p\right),\q \frac{r_0}{\sqrt\alpha}=\O\left((\chi h^{-1})^p\right).
\end{align}
\end{assumption}

\begin{theorem}
\label{thm:asymptotic-local-entry}
在\cref{eq:kappa-scaling,eq:schur-lambda-lower} 和\cref{ass:polynomial-data} 下, 存在
\begin{align}
    m_0&=\O\bigl(\sqrt\chi\,h^{-1}\log(\chi h^{-1})\bigr)
    \label{eq:local-entry-scale}
\end{align}
使得 PCG 要么在 $m_0$ 前精确终止, 要么对所有 $m\ge m_0$ 的后续非平凡迭代均有 $\norm{\bm Z_m}_{\mathrm F}\le \dt/8$. 固定 $\chi$ 时,
\begin{align}
    m_0&=\O(h^{-1}\log h^{-1})=o(h^{-2}).
    \label{eq:fixed-contrast-entry-scale}
\end{align}
\end{theorem}

\begin{proof}
由全局 PCG 误差界,
\begin{align}
    \norm{\bm Z_m}_{\mathrm F}
    &\le\frac{2r_0}{\sqrt{\lambda_{\min}(\bm S)}}c_3^m.
    \label{eq:Z-global-bound-proof-use}
\end{align}
因此要满足\cref{eq:local-ball-condition} 的局部条件 $\norm{\bm Z_m}_{\mathrm F}\le\dt/8$, 只需
\begin{align}
    c_3^m
    &\le\frac{\dt\sqrt{\lambda_{\min}(\bm S)}}{16r_0}.
    \label{eq:entry-sufficient-cg-factor}
\end{align}
即
\begin{align}
    m&\ge\frac{
        \log\left(
            \dfrac{16r_0}{\dt\sqrt{\lambda_{\min}(\bm S)}}
        \right)
    }{-\log c_3}:=m_0,
    \label{eq:entry-explicit-m0}
\end{align}
其中
\begin{align}
    -\log c_3
    &=\log\frac{\sqrt\kappa+1}{\sqrt\kappa-1}
      =2\operatorname{arctanh}(\kappa^{-1/2})
      \ge\frac2{\sqrt\kappa},
    \label{eq:cg-factor-log-lower}
\end{align}
故
\begin{align}
    \frac1{-\log c_3}
    &\le\frac{\sqrt\kappa}{2}
     =\O(\sqrt\chi\,h^{-1}).
    \label{eq:cg-log-denominator-scaling}
\end{align}
另一方面, $\lambda_{\min}(\bm S)\gtrsim\alpha$, 所以
\begin{align}
    \log\left(
        \frac{16r_0}{\dt\sqrt{\lambda_{\min}(\bm S)}}
    \right)&=\O\left(\log\left(
        \frac{r_0}{\sqrt\alpha}\dt^{-1}
    \right)\right)
     =\O(\log(\chi h^{-1})),
    \label{eq:entry-log-numerator-scaling}
\end{align}
其中最后一步使用\cref{ass:polynomial-data}. 将\cref{eq:cg-log-denominator-scaling,eq:entry-log-numerator-scaling} 代入\cref{eq:entry-explicit-m0}, 得
\begin{align}
    m_0
    &=\O\bigl(\sqrt\chi\,h^{-1}\log(\chi h^{-1})\bigr).
    \label{eq:local-entry-scale-proof}
\end{align}
固定 $\chi$ 后即得到\cref{eq:fixed-contrast-entry-scale}.
\end{proof}

由于二维细点数满足 $n_{\mathrm F}(h)=\Theta(h^{-2})$, 故\cref{eq:fixed-contrast-entry-scale} 还给出
\begin{align}
    \frac{m_0}{n_{\mathrm F}(h)}=\O(h\log h^{-1})\to0.
\end{align}
因此, \cref{thm:endpoint-local-expansion,thm:pcg-two-sided-step-criterion} 虽然在粗空间坐标中是局部结果, 但沿 PCG 路径在一个相对于二维代数终止尺度为低阶的初始过渡段之后均可适用.

\end{document}
